\documentclass[11pt, a4paper]{article}

\usepackage[T1]{fontenc}
\usepackage[utf8]{inputenc}
\usepackage{tgpagella}
\usepackage[spanish, es-noshorthands]{babel}
\usepackage{amsmath, amssymb, bm}
\usepackage{booktabs}
\usepackage{tabularx}
\usepackage[table, xcdraw]{xcolor}
\usepackage[most]{tcolorbox}
\usepackage[margin=2.5cm]{geometry}
\usepackage{fancyhdr}

\pagestyle{fancy}
\fancyhf{}
\setlength{\headheight}{16pt}
\lhead{\textbf{UCB Sede Tarija} | Ing. Mecatrónica}
\chead{}
\rhead{IMT-121 Circuitos Electrónicos I | Solucionario S07-1}
\lfoot{Prof: M.Sc. Ing. Bernardo Quiroga Turdera}
\rfoot{Página \thepage}
\renewcommand{\headrulewidth}{0.4pt}

\definecolor{ucbblue}{RGB}{0, 51, 102}

\begin{document}

\begin{tcolorbox}[colback=ucbblue!5!white, colframe=ucbblue, title=\textbf{Clave de Evaluación Docente: Experiencia Integrada A y Cargas}]
\end{tcolorbox}

\section*{Solucionario Artefacto 2 (Ejercicio Guiado de Barrido)}
Dato: $V_{th} = 12\,\text{V}$, $R_{th} = 3\,\text{k}\Omega$. Usando fórmulas $I_L = V_{th}/(R_{th}+R_L)$, $V_L = I_L R_L$, $P_L = I_L V_L$, $\eta = R_L / (R_L + R_{th})$.

\vspace{0.2cm}
\noindent\renewcommand{\arraystretch}{1.3}
\begin{tabularx}{\textwidth}{|>{\centering\arraybackslash}p{2cm}|X|X|X|X|X|}
\hline
\rowcolor{gray!10}
\textbf{$R_L$ (k$\Omega$)} & \textbf{$I_L$ (mA)} & \textbf{$V_L$ (V)} & \textbf{$P_L$ (mW)} & \textbf{$P_{Rth}$ (mW)} & \textbf{$\eta$ (\%)} \\
\hline
$0.75$ & $3.200$ & $2.400$ & $7.680$ & $30.720$ & $20\%$ \\
\hline
$1.5$  & $2.667$ & $4.000$ & $10.667$ & $21.336$ & $33.3\%$ \\
\hline
$\mathbf{3.0}$ & $\mathbf{2.000}$ & $\mathbf{6.000}$ & $\mathbf{12.000}$ & $\mathbf{12.000}$ & $\mathbf{50\%}$ \\
\hline
$6.0$  & $1.333$ & $8.000$ & $10.667$ & $5.333$ & $66.7\%$ \\
\hline
$12.0$ & $0.800$ & $9.600$ & $7.680$ & $1.920$ & $80\%$ \\
\hline
\end{tabularx}
\vspace{0.2cm}
\begin{itemize}\setlength\itemsep{-0.1em}
    \item Mayor $I_L$: En $0.75\,\text{k}\Omega$ (cargas que tienden al cortocircuito).
    \item Mayor $V_L$: En $12.0\,\text{k}\Omega$ (cargas que tienden al circuito abierto).
    \item Mayor $P_L$: En $3.0\,\text{k}\Omega$, donde $R_L = R_{th}$ (Punto de Máxima Transferencia). $P_{\text{max}} = 12^2/(4\cdot3000) = 12\,\text{mW}$.
    \item Mayor $\eta$: En $12.0\,\text{k}\Omega$ ($80\%$). La eficiencia sube monótonamente al incrementar $R_L$.
\end{itemize}

\section*{Solucionario Artefacto 3 (Prelaboratorio)}
\textbf{Parte 1.1: Superposición para $\bm{V_{oc}}$}\\
Actuando $V_s$ sola ($I_s$ abierto): El divisor en $a$ arroja $V_a^{(V_s)} = 12\,\text{V} \cdot \frac{8\,\text{k}}{4\,\text{k}+8\,\text{k}} = 8\,\text{V}$.\\
Actuando $I_s$ sola ($V_s$ en corto): Las resistencias quedan en paralelo $4\,\text{k}\Omega \parallel 8\,\text{k}\Omega = 2.667\,\text{k}\Omega$. \\
La contribución de voltaje es $V_a^{(I_s)} = 1\,\text{mA} \cdot 2.667\,\text{k}\Omega = 2.667\,\text{V}$.\\
Total $\bm{V_{oc} = V_{th}} = 8\,\text{V} + 2.667\,\text{V} = \mathbf{10.667\,\text{V}}$.

\vspace{0.2cm}
\noindent\textbf{Parte 1.2: Cálculo de $\bm{R_{th}}$}\\
Desactivando las fuentes ($V_s$ corto, $I_s$ abierto), miramos desde el puerto $a-b$:\\
$\mathbf{R_{th}} = 4\,\text{k}\Omega \parallel 8\,\text{k}\Omega = \mathbf{2.667\,\text{k}\Omega}$.

\vspace{0.2cm}
\noindent\textbf{Parte 1.3: Predicción Norton}\\
Por ley de conversión: $\mathbf{I_{sc}} = \mathbf{I_N} = \frac{10.667\,\text{V}}{2.667\,\text{k}\Omega} = \mathbf{4\,\text{mA}}$.

\vspace{0.2cm}
\noindent\textbf{Parte 1.4: Respuesta en Carga ($R_L = 8\,\text{k}\Omega$)}\\
Conectado a la serie de Thévenin: 
$\mathbf{I_L} = \frac{10.667\,\text{V}}{2.667\,\text{k}\Omega + 8\,\text{k}\Omega} = \mathbf{1\,\text{mA}}$. $\quad\mathbf{V_L} = 1\,\text{mA} \cdot 8\,\text{k}\Omega = \mathbf{8\,\text{V}}$.

\vspace{0.2cm}
\textbf{Nota de Seguridad Docente ($I_{sc}$):}\\
No solicite medir la corriente de cortocircuito directa de la placa/montaje si no ha validado que las resistencias térmicas limitadoras de la fuente o fusibles del multímetro aguantan la salida. En este circuito nominalmente son $4\,\text{mA}$, lo cual es inocuo para instrumentos estándar de lab, pero siempre evalúe las prácticas en función del hardware local antes de omitir la restricción.

\end{document}
